\chapter{Clase 2}

\section{Descripción de Datos y Análisis de Muestras}

\subsection{Técnicas de la Estadística Descriptiva}
Los estudios estadísticos, para poseer una cierta fiabilidad, deben basarse en la recogida de muchos datos: cuantos más, mejor. La Estadística dispone de técnicas específicas para cada etapa del tratamiento de la información:

\begin{longtblr}{colspec={|Q[l,wd=4cm]|Q[l,wd=9cm]|}, rowhead=1}
\hline
\textbf{Etapa} & \textbf{Descripción} \\
\hline
Recogida de datos & Obtener datos lo más representativos posible de la población. \\
\hline
Resumen en tablas & Organizar los datos en tablas ordenadas de frecuencias. \\
\hline
Representación gráfica & Representar gráficamente las tablas obtenidas. \\
\hline
Medidas representativas & Efectuar medidas centrales, de dispersión, asimetría, aplastamiento y correlaciones. \\
\hline
\end{longtblr}

Al estudio de estos aspectos se le denomina Estadística Descriptiva. La mayoría de los trabajos de nivel introductorio se limitan a estos aspectos.

\subsection{Inferencia a partir de Muestras}
Una vez seleccionada la muestra mediante los procedimientos estudiados en la Clase 1, sus resultados pueden utilizarse para estimar parámetros o contrastar afirmaciones acerca de la población. Este proceso de generalización constituye la Estadística Inferencial y exige cuantificar la incertidumbre asociada a las conclusiones.

\section{Contraste de Hipótesis}

En la actualidad, el contraste de hipótesis es la parte de la Estadística más usada en todo tipo de investigaciones. Ningún trabajo de nivel universitario o profesional se admite sin estar basado en un Diseño de Experimentos y en el uso de las técnicas de Contraste de Hipótesis.

\subsection{Concepto y Metodología}
Consiste en plantear una hipótesis (llamada nula) frente a otra alternativa, recoger datos representativos y comprobar si estos son consistentes con las hipótesis o no. Requiere de una metodología rigurosa.

Preguntas típicas que se resuelven mediante contraste de hipótesis:
\begin{itemize}
    \item ¿Ha mejorado el rendimiento de los estudiantes con el nuevo modelo de estudios?
    \item ¿Es el nuevo medicamento mejor que el anterior?
    \item ¿La duración de una bombilla es superior a mil horas?
\end{itemize}

Casi ningún producto entra en el mercado sin estudios estadísticos que apoyen la hipótesis con que se promocionan. Las técnicas usadas forman un cuerpo de teoría muy amplio llamado Estadística Inferencial y Diseño de Experimentos.

\subsection{Complemento: Tipos de Error en el Contraste}
Al realizar un contraste de hipótesis, existe la posibilidad de cometer dos tipos de error:

\begin{longtblr}{colspec={|Q[l,wd=3cm]|Q[l,wd=5cm]|Q[l,wd=5cm]|}, rowhead=1}
\hline
& \textbf{$H_0$ verdadera} & \textbf{$H_0$ falsa} \\
\hline
Rechazar $H_0$ & Error tipo I (probabilidad $\alpha$) & Decisión correcta \\
\hline
No rechazar $H_0$ & Decisión correcta & Error tipo II (probabilidad $\beta$) \\
\hline
\end{longtblr}

\section{Medida de Relaciones}

En la Sociedad y la Naturaleza se pueden descubrir relaciones y paralelismos que, en algunos casos, permiten representar mediante una fórmula una relación entre dos o más variables. Se llama Correlación al estudio de estas relaciones.

Ejemplos:
\begin{itemize}
    \item En la Sociedad: ¿El nivel de estudios está correlacionado con los ingresos?
    \item En la Naturaleza: ¿Las lluvias están correlacionadas con las cosechas?
\end{itemize}

Aunque el caso más frecuente es el de la comparación de dos variables, cada día se emprenden más estudios en los que algunas variables se relacionan con varias otras, llamadas explicativas.

Estos paralelismos no se deben confundir con relaciones causa-efecto, que tienen tratamientos más profundos. La correlación no implica causalidad.

\subsection{Complemento: Coeficiente de Correlación Lineal}
La fuerza de la asociación lineal entre dos variables se mide mediante el coeficiente de correlación lineal $r$, el cual satisface:

$$-1 \leq r \leq 1$$

Un valor de $r$ cercano a $+1$ indica una correlación lineal positiva fuerte; un valor cercano a $-1$ indica una correlación lineal negativa fuerte; un valor cercano a $0$ indica ausencia de correlación lineal.

\section{Establecimiento de Predicciones}

Los datos recogidos en un estudio pueden presentar tendencias o ciclos que quizás permitan predecir qué va a ocurrir fuera del rango de datos obtenido. Se distinguen dos tipos de predicción:

\begin{itemize}
    \item \textbf{Extrapolación:} predecir qué va a ocurrir en el futuro, fuera del rango observado.
    \item \textbf{Interpolación:} predecir valores intermedios dentro del rango observado.
\end{itemize}

Esta técnica también puede servir para completar datos perdidos o erróneos. En el establecimiento de las predicciones es fundamental conocer los márgenes de error. Las técnicas consiguientes están recogidas en la Teoría de la Regresión.

\section{El Método Estadístico}

El método estadístico es el procedimiento mediante el cual se sistematiza y organiza el proceso de aprendizaje iterativo para convertir los datos en información y esta en conocimiento:

$$\text{Datos} \xrightarrow{\text{Estadística Descriptiva}} \text{Información} \xrightarrow{\text{Inferencia Estadística}} \text{Conocimiento}$$

\subsection{El Método Científico}
El método científico es un método o procedimiento que ha caracterizado a la ciencia natural desde el siglo XVII, que consiste en la observación sistemática, medición, experimentación, la formulación, análisis y modificación de las hipótesis. Sus etapas son:

\begin{enumerate}
    \item Planteamiento del problema.
    \item Revisión de la literatura.
    \item Formulación de hipótesis.
    \item Observación y experimentación.
    \item Elaboración de conclusiones.
\end{enumerate}

\subsection{Estadística y Método Científico}
El método científico está compuesto por los siguientes pasos:
\begin{enumerate}
    \item Formular una teoría (problema).
    \item Recoger datos para probar la teoría.
    \item Analizar los datos.
    \item Interpretar los resultados y tomar una decisión.
\end{enumerate}

La ciencia está, por lo tanto, todo el tiempo revisando sus teorías. La Estadística no es un conjunto de técnicas aisladas unas de otras, sino que, en conjunto con el método científico, entrega un procedimiento analítico para tomar decisiones.

\subsection{Etapas del Método Estadístico}
Para tener éxito en una investigación estadística se deben seguir las siguientes etapas:

\begin{longtblr}{colspec={|Q[l,wd=3.5cm]|Q[l,wd=9.5cm]|}, rowhead=1}
\hline
\textbf{Etapa} & \textbf{Descripción} \\
\hline
Formulación del problema & Especificación del problema y definición de los objetivos. Crear conceptos precisos, formular preguntas claras, imponer límites teniendo en cuenta dinero y tiempo disponible. \\
\hline
Diseño del experimento & Obtener un máximo de información con un mínimo de error, dinero y tiempo. Determinar el tamaño de la muestra. Saber si la investigación es controlada u observacional. \\
\hline
Colección de datos & Muestreo: método usado para la recolección de datos. \\
\hline
Presentación y análisis & Ordenar datos en forma legible y condensada (diagramas, gráficas). Calcular medidas descriptivas básicas. \\
\hline
Inferencia y decisiones & Formular y obtener conclusiones o establecer generalizaciones acerca de la población. \\
\hline
\end{longtblr}

En resumen, se necesita:
\begin{enumerate}
    \item Definición del problema de estudio y objetivos del mismo.
    \item Selección de la información necesaria para realizar el estudio.
    \item Recogida de la información (depende del presupuesto y la calidad exigida).
    \item Ordenación y clasificación de la información en tablas y gráficos.
    \item Resumen de los datos mediante medidas de posición, dispersión, asimetría y concentración.
    \item Análisis estadístico formal obteniendo hipótesis y contrastándolas.
    \item Interpretación de resultados y extracción de conclusiones.
    \item Extrapolación y predicción.
\end{enumerate}

\section{Método Científico frente a Método Estadístico}

\begin{longtblr}{colspec={|Q[l,wd=6.5cm]|Q[l,wd=6.5cm]|}, rowhead=1}
\hline
\textbf{Método Científico} & \textbf{Método Estadístico} \\
\hline
Conjunto sistemático de procesos en los que se basa la ciencia para explicar cualquier fenómeno y las leyes que lo administran. & Adaptación del método científico a un área del conocimiento, aportando la planificación de la investigación. \\
\hline
Conjunto de procedimientos e instrumentos de la investigación científica que construye el camino reflexivo, sistemático, controlado y crítico para la adquisición del conocimiento. & Proceso de obtención, representación, simplificación, análisis, interpretación y proyección de las características de un estudio para una mejor comprensión de la realidad y una optimización en la toma de decisiones. \\
\hline
Sigue las interrogantes: observación, preguntas, hipótesis, experimentación, conclusiones, documentación y nuevos descubrimientos. & Herramienta poderosa de precisión científica en la medida en que se combine con los métodos cualitativos y se emplee de acuerdo a las necesidades y al sano criterio. \\
\hline
Consiste en la observación sistemática, medición y experimentación, y la formulación, análisis y modificación de la hipótesis. & Proporciona las técnicas necesarias para recolectar y analizar la información requerida. Se distingue una fase de planificación y otra de ejecución. \\
\hline
Ayuda a organizar adecuadamente la información de los fenómenos y a determinar las leyes que los rigen. & Comprende: recuento y compilación de datos, tabulación y agrupamiento, representación gráfica, medición de datos, inferencia estadística y predicción. \\
\hline
Estudio sistemático, controlado, empírico y crítico de proposiciones hipotéticas acerca de presuntas relaciones de varios fenómenos. & Secuencia de procedimientos para el manejo de los datos cualitativos y cuantitativos de la investigación. \\
\hline
\end{longtblr}
